\documentclass[11pt]{article}
\usepackage[margin=1in]{geometry}
\usepackage[T1]{fontenc}
\usepackage{xurl}
\usepackage[pdfencoding=auto]{hyperref}
\hypersetup{colorlinks=true,urlcolor=blue}
\usepackage{enumitem}
\setlength{\parskip}{4pt}
\begin{document}

\noindent
\textbf{To:} The Editor-in-Chief / Handling Editor, \textit{SoftwareX}\\[2pt]
\textbf{Date:} 15 October 2026\\[2pt]
\textbf{Article type:} Original Software Publication\\[2pt]
\textbf{Manuscript title:} \textit{Solderbench: An open-source synthetic benchmark for PCB solder-joint defect inspection algorithm development}\\[2pt]
\textbf{Manuscript number:} [insert manuscript number]\\[2pt]
\textbf{Revision:} Major Revision --- revised manuscript resubmitted

\vspace{8pt}
\noindent Dear Editor,

Thank you for the opportunity to revise our manuscript, and please pass our thanks to
both reviewers. Their comments exposed one genuinely serious problem in the original
submission, and addressing it has materially improved the paper. We have implemented
every requested change \textbf{in the manuscript and in the released software}, and we
resubmit a revised manuscript together with a point-by-point response.

\section*{What we changed --- summary}

\begin{enumerate}[leftmargin=*,itemsep=3pt]
\item \textbf{The S3DA baseline is now the original hand-crafted rule (the substantive fix).}
In the original submission our S3DA baseline was a shallow decision tree trained on all
seven features --- a learned proxy for, rather than a reproduction of, the hand-tuned
volume-threshold rule of Xu et al. [6]. We have now implemented the published rule itself
(\texttt{S3DARuleClassifier}: integral-summation volume compared against hand-set
acceptance limits of 0.90 and 1.10 times a nominal reference volume, the reference being
calibrated from training-set statistics only, with no access to test data or class labels)
and \textbf{re-run every experiment}. The reported S3DA accuracy changes accordingly, from
0.9131 $\pm$ 0.0156 (proxy) to \textbf{0.5632 $\pm$ 0.0091} (rule) over five seeds. All S3DA
numbers in the paper now come from the rule; the tree proxy is retained in the repository
for reference only and is used for no reported number. We additionally ship a
threshold-sensitivity sweep showing that the conclusion holds for every setting of the two
limits.

\item \textbf{All headline results are now multi-seed.} The single-seed comparison table is
replaced by a five-seed (42/123/456/789/1011) mean $\pm$ std table, and the abstract now
reports both the seed-42 reference value (0.9440 accuracy, 0.9439 macro-$F_1$) and the
five-seed statistic (\textbf{0.9531 $\pm$ 0.0115} for the random forest).

\item \textbf{Two new robustness studies.} A sensor-noise ablation ($\sigma$ from 0.005 to
0.10 mm) and a defect-strength / class-overlap ablation (narrow / default / wide) are
reported as new tables, with a figure showing the effect of noise on a nominal joint.

\item \textbf{The generator and the features are now fully specified.} A new table lists all
15 generator parameters with defaults and physical meaning (mm); a new figure shows one
representative point cloud per defect class; and the feature table now gives mathematical
definitions together with a computation/units column (voxel volume
$\mathit{vol} = N_{\mathrm{vox}}\delta^{3}$, $\delta = 0.01$ mm; mm and mm$^2$ units
throughout). We confirm that \textbf{$\delta = 0.01$ mm is used consistently} in the
manuscript and in the code --- no 0.05 mm value appears anywhere.

\item \textbf{Four over-claimed statements corrected.} (i) ``all learning-based methods
improve on the rule'' now reads ``most''; (ii) we no longer claim the random forest strictly
outperforms gradient boosting --- the mean gap (0.0053) is smaller than the inter-seed
standard deviation (0.0065), so the two are reported as statistically tied, with the term
defined; (iii) the SVM cross-validation/test gap is described as \emph{consistent with}
mild overfitting rather than as proof of it; (iv) excessive solder, previously called the
hardest defect, is corrected --- it is in fact well separated and incurs no test errors.

\item \textbf{Reproducibility and real-data scope.} Exact dependency versions are pinned
(Python 3.13; NumPy 2.5.2, SciPy 1.18.0, pandas 3.0.6, scikit-learn 1.9.0, Matplotlib
3.11.1), a \texttt{Dockerfile} builds a one-command reproducible image, the README has been
rewritten with parameters, results and reproduction commands with runtimes, and the test
suite now has 17 tests. Following Reviewer 1's comment 4, the manuscript now states
explicitly that \texttt{load\_dataset()} is a \textbf{software hook only}: real-world data
acquisition, joint segmentation and pre-processing are \textbf{not} implemented, and users
must supply pre-segmented solder-joint point clouds.

\item \textbf{A synthetic-only scope statement} has been added to the Limitations section:
the absolute accuracy values cannot be transferred directly to industrial inspection; only
the relative ordering under a controlled, identical protocol should be read as a result.
\end{enumerate}

\section*{Files enclosed}

\begin{itemize}[leftmargin=*,itemsep=2pt]
\item Revised manuscript \textbf{with track changes} (additions underlined/blue, deletions
struck out/red), file: \texttt{solderbench\_softwarex\_draft\_}\\ \texttt{v20\_trackchanges.pdf};
\item \texttt{solderbench\_softwarex\_draft\_v20\_clean.pdf} --- revised manuscript,
\textbf{clean copy} (identical content);
\item \texttt{RESPONSE\_TO\_REVIEWERS.md} --- point-by-point response to all ten reviewer comments.
\end{itemize}

\section*{Software and data availability}

The software accompanying this revision is release \textbf{v0.1.3}, available at
\url{https://github.com/wqinmin975/pcb} and archived on Zenodo under the concept DOI
\url{https://doi.org/10.5281/zenodo.22959192} (which always resolves to the latest
release). The benchmark dataset is not distributed as a fixed file: it is generated
deterministically at run time, and \texttt{python examples/run\_benchmark.py} regenerates
every table and figure in the paper.

\section*{Declarations}

\begin{itemize}[leftmargin=*,itemsep=2pt]
\item \textbf{Funding:} this research did not receive any specific grant from funding
agencies in the public, commercial, or not-for-profit sectors.
\item \textbf{Competing interests:} the authors declare none.
\item \textbf{CRediT authorship contribution statement:} included in the manuscript.
\item \textbf{Authors:} Qinmin Wang, Zixin Wang, Hongyuan Wang, Zongzhong Wang
(corresponding).
\item \textbf{Corresponding author:} Zongzhong Wang, School of Electromechanical
Engineering, Guangdong University of Technology, Guangzhou, China ---
\url{mailto:wangzz@gdut.edu.cn}.
\end{itemize}

We believe the revised manuscript and the accompanying software release fully address the
reviewers' concerns, and we thank you for your time and consideration.

\vspace{12pt}
\noindent Sincerely,\\[14pt]
\noindent \textbf{Zongzhong Wang} (on behalf of all authors)\\
School of Electromechanical Engineering, Guangdong University of Technology\\
wangzz@gdut.edu.cn

\end{document}
